% answers_short.tex —— 2012 年·湖北卷·理科 · 简短答案速查
% 由 tools/make_short_answers.py 自 answers.tex 的 \gksolhead/\gkanswer 提取，
% 并与 2012/2012_数学真题_湖北理_答案要点.md 逐题比对后生成，勿手改（改源 MD 后重跑）。

\begin{gknotice}
\textbf{使用说明}：本节只给最终答案便于快速核对，完整推导见解析版。答案为编写者独立推导，\textbf{非官方标准答案}。
\end{gknotice}

\gkpart{选择题}
\begin{center}\begin{tabular}{ccccccccccc}
\toprule
题号 & 1 & 2 & 3 & 4 & 5 & 6 & 7 & 8 & 9 & 10 \\
\midrule
答案 & A & D & B & B & D & C & C & A & C & D \\
\bottomrule
\end{tabular}\end{center}

\gkpart{填空题}
\begin{itemize}[leftmargin=2.6em]
  \item[\textbf{11.}] $\dfrac{2\pi}{3}$。条件化为 $a^2+b^2-c^2=-ab$，故 $\cos C=-\dfrac12$。
  \item[\textbf{12.}] $9$。循环依次给出 $(s,a,n)=(1,3,2)$、$(4,5,3)$、$(9,7,3)$，此时 $n<3$ 不成立，输出 $9$。
  \item[\textbf{13.}] (1) $90$。4 位回文数形如 $\overline{abba}$，首位 $9$ 种、次位 $10$ 种，共 $9\times10=90$ 个；(2) $9\times10^n$。$2n+1$ 位回文数由前 $n+1$ 位决定，首位 $9$ 种、其余 $n$ 位各 $10$ 种
  \item[\textbf{14.}] (1) $\dfrac{1+\sqrt5}{2}$。由圆心到菱形边的距离等于 $a$ 得 $b^2c^2=a^2(b^2+c^2)$，化为 $e^4-3e^2+1=0$；(2) $\dfrac{2+\sqrt5}{2}$。$S_1=2bc$，$S_2=\dfrac{4b^3c^3}{(b^2+c^2)^2}$，比值化简为 $\dfrac{2e^2-1}{2}$
  \item[\textbf{15.}] $2$。设圆心到弦的距离为 $d_0$，$D$ 到垂足的距离为 $t$，则 $CD^2=4-t^2\leqslant4$，当 $D$ 为 $AB$ 中点时取等。
  \item[\textbf{16.}] $\left(\dfrac52,\dfrac52\right)$。射线为 $y=x$（$x\geqslant0$），代入参数方程得 $t^2-3t=0$，交点 $(1,1)$、$(4,4)$，中点即所求。
\end{itemize}

\gkpart{解答题}
\begin{itemize}[leftmargin=2.6em]
  \item[\textbf{17.}] (1) $\dfrac{6\pi}{5}$。化简得 $f(x)=2\sin\left(2\omega x-\dfrac\pi6\right)+\lambda$，由对称轴 $x=\pi$ 与 $\omega\in\left(\dfrac12,1\right)$ 得 $\omega=\dfrac56$；(2) $\left[-1-\sqrt2,\,2-\sqrt2\right]$。由过点 $\left(\dfrac\pi4,0\right)$ 得 $\lambda=-\sqrt2$，再求 $2\sin\left(\dfrac53x-\dfrac\pi6\right)-\sqrt2$ 在 $\left[0,\dfrac{3\pi}5\right]$ 上的值域
  \item[\textbf{18.}] (1) $a_n=5-3n$ 或 $a_n=3n-7$。由 $3a_2=-3$ 得 $a_2=-1$，$a_1a_3=-8$，$a_1+a_3=-2$，公差为 $\pm3$；(2) $S_1=4$，$n\geqslant2$ 时 $S_n=\dfrac{3n^2-11n+20}{2}$。等比条件筛选出 $a_n=3n-7$，其前三项为 $-4,-1,2$
  \item[\textbf{19.}] (1) $BD=1$。此时 $DC=AD=2$，$V=\dfrac16x(3-x)^2$ 在 $x=1$ 处取最大值 $\dfrac23$；(2) $N$ 为棱 $CD$ 上满足 $DN=\dfrac12$（即 $DN=\dfrac14DC$）的点；$EN$ 与平面 $BMN$ 所成角为 $60^\circ$。以 $D$ 为原点建立坐标系，由 $\overrightarrow{EN}\cdot\overrightarrow{BM}=0$ 定出 $N$，再求线面角
  \item[\textbf{20.}] (1) $E(Y)=3$，$D(Y)=9.8$。$Y$ 的分布为 $P(Y=0)=0.3$，$P(Y=2)=0.4$，$P(Y=6)=0.2$，$P(Y=10)=0.1$；(2) $\dfrac67$。条件概率为 $\dfrac{P(300\leqslant X<900)}{P(X\geqslant300)}=\dfrac{0.6}{0.7}$
  \item[\textbf{21.}] (1) $C:\ x^2+\dfrac{y^2}{m^2}=1$，是椭圆；$m>1$ 时焦点为 $\left(0,\pm\sqrt{m^2-1}\right)$，$0<m<1$ 时焦点为 $\left(\pm\sqrt{1-m^2},0\right)$。由 $|DM|=m|DA|$ 得 $y=\pm m\sin\theta$，消去参数即得；(2) 存在，$m=\sqrt2$。把直线 $QN$ 与 $C$ 的第二交点 $H$ 用 $k$、$m$ 表示后，$PQ\perp PH$ 化为与 $k$ 无关的方程 $m^2=2$
  \item[\textbf{22.}] (1) $0$。$f'(x)=r\left(1-x^{r-1}\right)$ 在 $x=1$ 变号，$f(1)=0$ 为最小值；(2) 见解析版；(3) 推广为“设 $a_i\geqslant0$，$b_i$ 为正有理数（$i=1,2,\cdots,n$），$\sum\limits_{i=1}^n b_i=1$，则 $\prod\limits_{i=1}^n a_i^{b_i}\leqslant\sum\limits_{i=1}^n b_i a_i$”，对 $n$ 用数学归纳法证明，归纳步把前 $n$ 项的权归一化后先用归纳假设、再用 $n=2$ 的情形
\end{itemize}
